% =====================================================================
\chapter{État de l'art}
\label{chap-etat_de_lart}
% =====================================================================

Ce chapitre situe DOM-Graph RAG dans son contexte scientifique et se limite à ce qui est strictement nécessaire pour situer les choix de conception~:
le problème du découpage documentaire et ses métriques
(\cref{sec-eda_rag}), le prétraitement des documents visuellement riches
(\cref{sec-eda_docvisuel}), la famille du RAG visuel et son représentant le
plus proche, PixelRAG (\cref{sec-eda_ragvisuel}), le RAG structuré par graphe
et son ancêtre direct, MAGE-RAG (\cref{sec-eda_graphe}), et l'apprentissage
contrastif mobilisé par le lecteur (\cref{sec-eda_contrastif}). Il se clôt sur
le positionnement précis de ce travail (\cref{sec-eda_synthese}).

% =====================================================================
\section{Génération augmentée par récupération et découpage}
\label{sec-eda_rag}
% =====================================================================

La \cref{fig-eda_rag_chaine} donne les trois opérateurs qu'enchaîne un système
RAG. Le découpeur $C : \mathcal{D} \mapsto \mathcal{U}$ transforme une
collection de documents en unités récupérables. Le récupérateur
$R : q \mapsto \mathcal{U}_k \subseteq \mathcal{U}$ sélectionne les $k$ unités
les plus pertinentes pour une question $q$ selon un score $s(q,u)$. Le
générateur $G : (q, \mathcal{U}_k) \mapsto a$ répond à partir de la question
et de ces unités.

\begin{figure}[H]
\centering
\begin{tikzpicture}[
  every node/.style={font=\small},
  op/.style={draw, rounded corners=2pt, minimum width=1.75cm, minimum height=1.0cm,
             align=center, fill=codebg, draw=rulegray, line width=0.7pt},
  focus/.style={op, fill=accent!8, draw=accent, line width=1pt},
  arr/.style={-{Stealth[length=2.0mm]}, thick, draw=black!55},
  lbl/.style={font=\scriptsize, text=black!60, align=center}
]
\node (D) {$\mathcal{D}$};
\node[focus, right=0.5cm of D] (C) {$C$\\[-4pt]{\scriptsize découpeur}};
\node[right=0.5cm of C] (U) {$\mathcal{U}$};
\node[op, right=0.5cm of U] (R) {$R$\\[-4pt]{\scriptsize récupérateur}};
\node[right=0.5cm of R] (Uk) {$\mathcal{U}_k$};
\node[op, right=0.5cm of Uk] (G) {$G$\\[-4pt]{\scriptsize générateur}};
\node[right=0.5cm of G] (a) {$a$};

\draw[arr] (D) -- (C);
\draw[arr] (C) -- (U);
\draw[arr] (U) -- (R);
\draw[arr] (R) -- (Uk);
\draw[arr] (Uk) -- (G);
\draw[arr] (G) -- (a);

\node[lbl, below=1.15cm of Uk] (q) {question $q$};
\draw[arr] (q) -| (R.south);
\draw[arr] (q) -| (G.south);

\node[draw=black!35, dashed, rounded corners=3pt, inner sep=5pt,
      fit=(R)(Uk)(G)] (box) {};
\node[lbl, above=1pt of box.north] {optimisés par la littérature parcourue};
\node[lbl, above=0.30cm of C] {objet de ce travail};
\node[lbl, below=0.28cm of C, text width=2.9cm]
  {fixe le plafond mesuré en \cref{sec-eda_pool}};
\end{tikzpicture}
\caption[Les trois opérateurs d'un système RAG]{Les trois opérateurs d'un
système RAG. La composition impose un ordre de dépendance, car $R$ ne peut
sélectionner que des unités produites par $C$, et $G$ ne peut répondre qu'à
partir de ce que $R$ transmet.}
\label{fig-eda_rag_chaine}
\end{figure}
\citet{rag2020} ont formalisé cette architecture et montré qu'affiner
conjointement l'encodeur de requête et le générateur améliore les tâches à
forte intensité de connaissances~; \citet{fid2021} ont montré qu'encoder
chaque passage séparément avant de les fusionner dans le décodeur permet d'en
traiter un grand nombre sans coût quadratique. \citet{ragsurvey2023} recensent
ces variantes sans qu'aucune ne s'impose.

Le score $s(q,u)$ oppose deux familles. Les approches lexicales
(TF-IDF, BM25~\citep{bm25_2009}) ne demandent aucun entraînement et restent
compétitives~: le \anglais{benchmark} BEIR~\citep{beir2021} établit qu'en
transfert zéro-coup BM25 n'est pas systématiquement dépassé par les
récupérateurs denses, et la \cref{sec-res_ablation_b} retrouve ce constat.
Leur limite est le décalage de vocabulaire. Les approches denses
projettent question et unité dans un espace vectoriel commun appris par
contraste~\citep{sbert2019, simcse2021}~; \citet{dpr2020} en tirent la recette
de référence pour la récupération de passages, et ColBERT~\citep{colbert2020}
en conserve un vecteur par token plutôt qu'un vecteur unique, une idée qui
réapparaît dans le RAG visuel (\cref{sec-eda_ragvisuel}).

Cette littérature suppose $C$ donné, alors qu'il fixe un plafond que ni $R$ ni
$G$ ne peuvent franchir. Si la réponse à $q$ est répartie sur deux unités
distinctes, ou noyée dans une unité majoritairement hors sujet, aucun score ne
peut compenser. Les stratégies de découpage textuel, fenêtres de taille fixe, découpage guidé
par la structure du document ou découpage sémantique aux points de rupture de
similarité entre plongements~\citep{semchunk2025}, sont recensées par
\citet{ragsurvey2023} sans que leur effet sur ce plafond soit mesuré. Toutes
opèrent sur du texte déjà extrait, où les frontières de mise en page ont été
détruites, alors qu'un moteur de rendu les calcule sans coût additionnel sur le
web (\cref{sec-rendu_dom}). Toutes mesurent aussi la qualité du classement en
supposant l'attachement de la réponse acquis, ce que ne fait ni la
\cref{sec-eda_pool} ni la \cref{sec-res_a_attach}.

\subsection{Le biais de taille de vivier}
\label{sec-eda_pool}

Comparer deux stratégies de découpage produisant des viviers de candidats de
tailles différentes pose un piège méthodologique peu discuté~: sur une même
page en \num{5} bandes contre \num{25} tuiles, un classement aléatoire place
la bonne unité au premier rang avec une probabilité de \SI{20}{\percent} dans
le premier cas et de \SI{4}{\percent} dans le second, de sorte que comparer
des Recall@1 bruts revient en partie à comparer des tailles de vivier. Le rang
percentile

\begin{equation}
\mathrm{pct}(u^{+}) = 1 - \frac{\mathrm{rang}(u^{+}) - 1}{n_{\text{items}} - 1},
\label{eq-pctrank}
\end{equation}

vaut $0{,}5$ en espérance sous classement aléatoire quelle que soit la taille
du vivier, et constitue la métrique de comparaison entre conditions de
tuilage employée en \cref{sec-res_ablation_a}, aux côtés du taux
d'attachement qui représente la proportion de questions dont la réponse survit au
découpage, absente des protocoles usuels alors qu'elle détermine le plafond
que ni $R$ ni $G$ ne peuvent franchir.

% =====================================================================
\section{Documents visuellement riches et prétraitement}
\label{sec-eda_docvisuel}
% =====================================================================

Un document est dit visuellement riche lorsque sa mise en page porte de
l'information sémantique, comme un tableau, une facture ou une page à encadrés.
Pour un tel document, la chaîne de prétraitement classique enchaîne trois
étapes à perte, détaillées par la \cref{fig-eda_pertes}. Le résultat cumulé est
qu'une question du type \og quelle est la valeur de la colonne revenus pour la
ligne \num{2024}\fg{} devient difficile à satisfaire, non parce que le
récupérateur est mauvais, mais parce que la correspondance ligne-colonne a été
détruite en amont.

\begin{figure}[H]
\centering
\begin{tikzpicture}[
  every node/.style={font=\small},
  stage/.style={draw, rounded corners=2pt, minimum width=2.55cm, minimum height=0.95cm,
                align=center, fill=codebg, draw=rulegray, line width=0.7pt,
                font=\scriptsize},
  arr/.style={-{Stealth[length=2.0mm]}, thick, draw=black!55},
  etape/.style={font=\scriptsize\itshape, text=black!60, align=center},
  loss/.style={font=\scriptsize, text=black!60, align=center, text width=2.75cm}
]
\node[stage] (doc) {Document\\mis en page};
\node[stage, right=1.6cm of doc] (seq) {Séquence de\\caractères};
\node[stage, right=1.6cm of seq] (ord) {Texte\\réordonné};
\node[stage, right=1.6cm of ord] (unit) {Unités\\récupérables $\mathcal{U}$};

\draw[arr] (doc) -- node[etape, above, text width=1.7cm] {extraction\\ou OCR} (seq);
\draw[arr] (seq) -- node[etape, above, text width=1.7cm] {ordre de\\lecture} (ord);
\draw[arr] (ord) -- node[etape, above, text width=1.7cm] {découpage\\$C$} (unit);

\node[loss, anchor=north] at ($(doc.south)!0.5!(seq.south)+(0,-0.35)$)
  {la structure 2D devient linéaire~; l'OCR y ajoute des erreurs de caractères};
\node[loss, anchor=north] at ($(seq.south)!0.5!(ord.south)+(0,-0.35)$)
  {l'ordre du fichier n'est pas l'ordre de lecture~; heuristiques à taux
   d'erreur non nul};
\node[loss, anchor=north] at ($(ord.south)!0.5!(unit.south)+(0,-0.35)$)
  {frontières d'unité fixées sans égard aux frontières de contenu
   (\cref{sec-eda_rag})};
\end{tikzpicture}
\caption[Les trois pertes de la chaîne de prétraitement classique]{Les trois
étapes à perte de la chaîne de prétraitement classique et l'information détruite
à chacune. Les modèles sensibles à la mise en page réinjectent la géométrie en
aval de la première étape~; les modèles sans OCR suppriment les deux premières.}
\label{fig-eda_pertes}
\end{figure}

Une première réponse injecte la géométrie dans la représentation.
LayoutLM~\citep{layoutlm2020} ajoute aux plongements de tokens des positions 2D
dérivées de la boîte englobante de chaque mot~; LayoutLMv2~\citep{layoutlmv2_2021}
et LayoutLMv3~\citep{layoutlmv3_2022} l'intègrent au pré-entraînement (\num{92.1}
de $F_1$ sur FUNSD, \num{83.4} d'ANLS sur DocVQA~\citep{docvqa2021})~; MarkupLM~\citep{markuplm2022}
transpose l'idée au web en encodant le chemin XPath de chaque token. Ces
modèles reçoivent tous des mots et boîtes déjà produits par un OCR en amont, et
la structure y est apprise, donc coûteuse à ré-entraîner sur un type de
document non couvert~; surtout, elle sert la représentation d'une unité déjà
découpée, non la décision de découpage~: MarkupLM sait que deux tokens
partagent un ancêtre DOM, il ne décide pas où couper la page.

Une seconde réponse supprime l'OCR~: Donut~\citep{donut2022} produit du texte
structuré directement depuis l'image d'un document, et
Pix2Struct~\citep{pix2struct2023} porte l'idée au web en pré-entraînant sur
\num{80} millions de captures d'écran, avec pour cible un HTML simplifié dérivé
du DOM rendu. Cela permet ainsi d'établir que le DOM porte la structure dont un
modèle visuel a besoin, mais il l'exploite comme cible de supervision apprise,
quand ce travail l'emploie comme décision de découpage à l'inférence, sans
apprentissage.

% =====================================================================
\section{Modèles vision-langage et RAG visuel}
\label{sec-eda_ragvisuel}
% =====================================================================

Les modèles vision-langage contemporains combinent une tour de vision de type
\anglais{Vision Transformer}~\citep{vit2021}, qui découpe l'image en
\anglais{patches} traités par auto-attention~\citep{transformer2017}, et un
modèle de langue décodeur consommant une séquence mixte de tokens visuels et
textuels, sur la base des représentations image-texte alignées par
contraste~\citep{clip2021}. Qwen2-VL~\citep{qwen2vl2024}, employé au
\cref{chap-methodologie}, apporte le support de la résolution dynamique~: le
nombre de tokens visuels y est proportionnel à l'aire réelle de l'image, une
propriété dont le \cref{sec-eda_cout_visuel} tire les conséquences pour ce
travail.

Si un tel modèle produit un plongement d'image aligné sur le langage, on peut
indexer des images de pages entières et récupérer dessus~: \anglais{Document
Screenshot Embedding}~\citep{dse2024} le fait sans extraction de contenu
préalable, et ColPali~\citep{colpali2024} porte le mécanisme d'interaction
tardive de ColBERT aux \anglais{patches} visuels.

\subsection{PixelRAG et le découpage aveugle}
\label{sec-eda_pixelrag}

Le travail le plus directement relié est PixelRAG~\citep{pixelrag2026}, qui
applique cette approche à l'échelle du web~: les pages sont rendues, découpées
en bandes et indexées avec un encodeur vision-langage, et les auteurs
rapportent égaler ou dépasser la récupération textuelle sur plusieurs bancs
d'essai, avec une réduction du coût en tokens allant jusqu'à un facteur trois
par sous-échantillonnage de la résolution. C'est son découpage qui nous
intéresse~: une grille de hauteur fixe, la page rendue tranchée en bandes de
$\num{875} \times \num{1024}$ pixels indépendamment de leur contenu. La
décision est simple et sans dépendance au type de document, mais son défaut
est celui de la fenêtre de taille fixe qui crée une bande coupant là où elle tombe,
au milieu d'un tableau comme d'un paragraphe. C'est cette condition qui est
réimplémentée (\cref{sec-protocole_grille}) et mesurée
(\cref{sec-res_ablation_a}).

% =====================================================================
\section{RAG structuré par graphe}
\label{sec-eda_graphe}
% =====================================================================

Traiter $\mathcal{U}$ comme un sac plat empêche de récupérer une unité peu
pertinente lexicalement mais indispensable contextuellement, et interdit le
raisonnement multi-sauts. Deux issues coexistent~: itérer la récupération ou
structurer les preuves en amont. La première voie est celle de
ReAct~\citep{react2023}, IRCoT~\citep{ircot2023}, FLARE~\citep{flare2023} et
Self-RAG~\citep{selfrag2024}, qui entrelacent raisonnement et appels de
récupération, au prix d'appels au modèle en série. Dans la seconde,
\citet{graphrag2024} construisent un graphe de connaissances
sur des entités extraites par un modèle de langue, ce qui répond à des
questions globales sur un corpus qu'aucune récupération locale ne
satisferait~; le graphe porte sur des entités, non sur la structure
documentaire.

\subsection{MAGE-RAG}
\label{sec-eda_magerag}

\citet{magerag2026} construisent au contraire un graphe de la structure
du document, et c'est l'ancêtre direct de l'architecture proposée ici~: des
nœuds de niveau page et de niveau élément, reliés par cinq relations
(\code{contains}, \code{reading\_order}, \code{layout\_adjacency},
\code{section\_hierarchy}, \code{semantic\_neighbor}), que parcourt à la
requête un contrôleur budgété décidant quels nœuds le modèle lecteur ouvre
réellement. L'\cref{sec-controleur} en généralise l'algorithme. La
différence tient à la provenance du graphe~: MAGE-RAG le construit hors ligne,
à partir d'un document déjà segmenté et d'un ordre de lecture déjà
déterminé, quand celui construit ici l'est depuis le DOM d'une page rendue, qui
fournit sans coût additionnel frontières exactes et types d'éléments
(\cref{sec-rendu_dom}). Deux relations que le contexte web rend
nécessaires s'y ajoutent, \code{links\_to} et \code{continues}
(\cref{tab-graphe_relations}). MAGE-RAG est évalué sur LongDocURL et
MMLongBench-Doc, deux jeux de données constitués de PDF déjà analysés~; ces
chiffres sont rapportés au \cref{chap-resultats} comme référence de
l'architecture, non comme comparaison numérique, faute d'un DOM à exploiter
sur un PDF.

% =====================================================================
\section{Apprentissage contrastif et adaptation efficace}
\label{sec-eda_contrastif}
% =====================================================================

\subsection{InfoNCE et la qualité des négatifs}
\label{sec-eda_contrastif_infonce}

Entraîner un récupérateur revient à apprendre une similarité $s(q,u)$ qui
place la bonne unité plus près de la question que les mauvaises. L'objectif
standard, InfoNCE~\citep{cpc2018}, largement repris en vision~\citep{simclr2020},
s'écrit pour une question de plongement $\mathbf{z}_q$, une unité positive
$u^+$ et des négatifs $\{u_1^-,\dots,u_K^-\}$~:

\begin{equation}
\mathcal{L}_{\text{InfoNCE}}
= -\log
\frac{\exp\!\big(\mathrm{sim}(\mathbf{z}_q, \mathbf{z}_{u^{+}}) / \tau\big)}
     {\exp\!\big(\mathrm{sim}(\mathbf{z}_q, \mathbf{z}_{u^{+}}) / \tau\big)
      + \sum_{i=1}^{K}
      \exp\!\big(\mathrm{sim}(\mathbf{z}_q, \mathbf{z}_{u^{-}_i}) / \tau\big)},
\label{eq-infonce}
\end{equation}

où $\tau$ est une température qui concentre le gradient sur les négatifs
proches du positif à mesure qu'elle diminue. Tout dépend des négatifs~: les
\anglais{in-batch negatives}~\citep{dpr2020} réutilisent gratuitement les
positifs des autres exemples du lot, tandis que les négatifs difficiles sont
minés délibérément. Le \cref{sec-contrastif} reprend cette équation, l'étend
au lot entier, et l'analyse plus en détail. La génération synthétique de
paires question-passage, établie par InPars~\citep{inpars2022} et
Promptagator~\citep{promptagator2022}, est transposée ici à des tuiles
d'image.

\subsection{Adaptation efficace d'un modèle génératif en encodeur}
\label{sec-eda_contrastif_adaptation}

Affiner intégralement un modèle de plusieurs milliards de paramètres est hors
de portée d'une infrastructure académique modeste. LoRA~\citep{lora2021} gèle
les poids pré-entraînés et n'apprend, pour chaque matrice ciblée, qu'une mise à
jour de rang faible~; QLoRA~\citep{qlora2023} y ajoute la quantification du
modèle gelé. Dans le cas multimodal, LoRA peut cibler le décodeur seul ou
aussi la tour de vision puisque la tâche traitée au \cref{sec-contrastif_lora}
exige les deux. Un second obstacle est que ces modèles sont génératifs et
n'exposent aucun vecteur de représentation~; PromptEOL~\citep{prompteol2024}
fait suivre l'entrée d'une amorce à compléter et prend l'état caché du
dernier token comme plongement, un procédé appliqué à l'identique à une
question textuelle et à une image de tuile au \cref{sec-contrastif_encodeur}.

% =====================================================================
\section{Synthèse et positionnement}
\label{sec-eda_synthese}
% =====================================================================

Trois constats délimitent une lacune précise, que le
\cref{tab-eda_positionnement} résume.

\begin{table}[H]
\centering
\small
\caption[Positionnement de DOM-Graph RAG]{Origine des frontières de l'unité
récupérable chez les travaux les plus proches. Elles sont soit arbitraires,
fenêtre ou grille, soit héritées d'un découpage effectué ailleurs, format du
fichier source ou corpus déjà segmenté en amont.}
\label{tab-eda_positionnement}
\begin{tabular}{@{}ll@{}}
\toprule
\textbf{Travail} & \textbf{Origine des frontières} \\
\midrule
\multicolumn{2}{@{}l}{\emph{Récupération textuelle}} \\
\quad DPR~\citep{dpr2020} & fenêtre fixe \\
\quad GraphRAG~\citep{graphrag2024} & fenêtre fixe \\[3pt]
\multicolumn{2}{@{}l}{\emph{Modèles sensibles à la mise en page}} \\
\quad LayoutLM~\citep{layoutlm2020} & corpus segmenté en amont \\
\quad MarkupLM~\citep{markuplm2022} & corpus segmenté en amont \\[3pt]
\multicolumn{2}{@{}l}{\emph{RAG visuel, lecture des pixels}} \\
\quad DSE~\citep{dse2024} & format du fichier source, la page \\
\quad ColPali~\citep{colpali2024} & format du fichier source, la page \\
\quad PixelRAG~\citep{pixelrag2026} & grille de hauteur fixe \\[3pt]
\multicolumn{2}{@{}l}{\emph{RAG par graphe de structure}} \\
\quad MAGE-RAG~\citep{magerag2026} & corpus segmenté en amont \\
\midrule
\textbf{Ce travail}, pixels et graphe & \textbf{boîtes englobantes du DOM rendu} \\
\bottomrule
\end{tabular}
\end{table}
Premièrement, le RAG visuel démontre que lire les pixels préserve une
information que l'extraction textuelle détruit, mais ne traite pas le
découpage~: soit il l'évite en travaillant à l'échelle de la page, soit il le tranche à
l'aveugle, comme PixelRAG. Deuxièmement, les modèles sensibles à la
mise en page exploitent la structure, mais en l'apprenant à grands frais et
pour représenter des unités déjà découpées. Enfin, MAGE-RAG
démontre qu'un contrôleur de graphe budgété surpasse une sélection top-$k$
statique, mais suppose les éléments du document déjà constitués en amont de
son propre pipeline, c'est-à-dire l'artefact que PixelRAG jette au rendu et que
ce travail restitue.

La thèse défendue est que, pour une page web sémantiquement balisée, la
segmentation dont le RAG visuel a besoin est déjà disponible et exacte dans le
DOM de la page rendue, et qu'il est donc inutile de la reconstruire depuis les
pixels ou de la contourner par une grille. Le \cref{chap-methodologie} formalise
et met à l'épreuve cette thèse, en commençant précisément là où PixelRAG
s'arrête~: par ce qu'un navigateur expose de la géométrie d'une page rendue.